\documentclass[11pt,a4paper]{article}
% =====================================================================
%  Shared preamble for the Part V lecture notes (lectures 19-22).
%
%  These four lectures sit outside Geron, so the notes ARE the primary
%  source and are examined on the same terms as the textbook parts.
%  Everything here is chosen to make that readable on paper: one column,
%  generous measure, and the same six-colour palette as the slides so a
%  student moving between the two is not re-learning what red means.
%
%  Build with pdflatex (twice, for the toc and the refs):
%      cd notes && latexmk -pdf lecture-19.tex
%  or  make -C notes
% =====================================================================
\usepackage[T1]{fontenc}
\usepackage[utf8]{inputenc}
\usepackage{newtxtext,newtxmath}      % Times-like: prints well, reads well
\usepackage[scaled=0.86]{inconsolata} % code, at a size that sits beside Times
\usepackage{microtype}
% newtx defines \Bbbk, \openbox, \proof and \endproof; amssymb and amsthm
% then redefine them and abort. Freeing the four names before those two
% packages load is the standard fix, and it costs nothing here: we use
% amsthm only for \newtheorem.
\let\Bbbk\relax
\usepackage{amsmath,amssymb}
\let\openbox\relax
\let\proof\relax
\let\endproof\relax
\usepackage{amsthm}
\usepackage{booktabs}
\usepackage{enumitem}
\usepackage{xcolor}
\usepackage{tcolorbox}
\usepackage{titlesec}
\usepackage{graphicx}
\usepackage[hidelinks]{hyperref}
\tcbuselibrary{skins,breakable}

% ---- the course palette, from assets/css/custom.css -------------------
\definecolor{aimlprimary}{HTML}{0B3D62}   % deep blue  - structure
\definecolor{aimlaccent} {HTML}{C0392B}   % brick red  - failures
\definecolor{aimlsuccess}{HTML}{14663A}   % green      - repairs
\definecolor{aimlmath}   {HTML}{6C3483}   % purple     - the derivation
\definecolor{aimlmuted}  {HTML}{4B5563}
\definecolor{aimlrule}   {HTML}{B0BCC7}
\definecolor{aimlsoft}   {HTML}{F4F7F9}

\usepackage[a4paper,margin=32mm,top=30mm,bottom=30mm]{geometry}
\setlength{\parskip}{0.45\baselineskip}
\setlength{\parindent}{0pt}
\linespread{1.06}

% ---- headings --------------------------------------------------------
\titleformat{\section}{\Large\bfseries\color{aimlprimary}}{\thesection}{0.7em}{}
\titleformat{\subsection}{\large\bfseries\color{aimlprimary}}{\thesubsection}{0.7em}{}
\titleformat{\subsubsection}{\bfseries\color{aimlmuted}}{}{0em}{}
\titlespacing*{\section}{0pt}{1.6\baselineskip}{0.5\baselineskip}
\titlespacing*{\subsection}{0pt}{1.2\baselineskip}{0.35\baselineskip}

% ---- boxes -----------------------------------------------------------
% One box per role, and the role is the colour. A student who has seen the
% slides already knows what each one means before reading a word of it.
\newtcolorbox{keybox}[1][]{
  breakable, enhanced, colback=aimlsoft, colframe=aimlprimary,
  boxrule=0.4pt, left=8pt, right=8pt, top=6pt, bottom=6pt,
  fonttitle=\bfseries\small, coltitle=white, colbacktitle=aimlprimary,
  attach boxed title to top left={yshift=-2mm, xshift=4mm},
  boxed title style={boxrule=0pt, arc=1pt}, #1}

\newtcolorbox{failbox}[1][]{
  breakable, enhanced, colback=white, colframe=aimlaccent,
  boxrule=0.9pt, left=8pt, right=8pt, top=6pt, bottom=6pt,
  fonttitle=\bfseries\small, coltitle=white, colbacktitle=aimlaccent,
  attach boxed title to top left={yshift=-2mm, xshift=4mm},
  boxed title style={boxrule=0pt, arc=1pt}, #1}

\newtcolorbox{derivbox}[1][]{
  breakable, enhanced, colback=white, colframe=aimlmath,
  boxrule=0.6pt, left=8pt, right=8pt, top=6pt, bottom=6pt,
  fonttitle=\bfseries\small, coltitle=white, colbacktitle=aimlmath,
  attach boxed title to top left={yshift=-2mm, xshift=4mm},
  boxed title style={boxrule=0pt, arc=1pt}, #1}

% An exercise. Deliberately the same shape as the notebook's `try` field:
% one change, and what should happen to the output.
\newtcolorbox{trybox}[1][]{
  breakable, enhanced, colback=white, colframe=aimlsuccess,
  boxrule=0.5pt, left=8pt, right=8pt, top=6pt, bottom=6pt,
  fonttitle=\bfseries\small, coltitle=white, colbacktitle=aimlsuccess,
  attach boxed title to top left={yshift=-2mm, xshift=4mm},
  boxed title style={boxrule=0pt, arc=1pt}, #1}

% ---- small semantics -------------------------------------------------
\newcommand{\scope}[1]{\textcolor{aimlmuted}{\small #1}}
\newcommand{\term}[1]{\textbf{#1}}
\newcommand{\code}[1]{\texttt{\small #1}}
\newcommand{\lecture}[1]{Lecture~#1}

\newtheorem{definition}{Definition}[section]
\newtheorem{proposition}[definition]{Proposition}

% ---- title block -----------------------------------------------------
% \notestitle{19}{Information retrieval: the lexical foundation}{SciFact (BEIR)}
\newcommand{\notestitle}[3]{%
  \thispagestyle{empty}
  \begin{flushleft}
    {\color{aimlmuted}\small\sffamily
     APPLICATIONS OF MACHINE LEARNING \quad$\cdot$\quad
     BSc Mathematics of Artificial Intelligence}\\[2mm]
    {\color{aimlrule}\rule{\textwidth}{0.8pt}}\\[4mm]
    {\color{aimlmuted}\large Lecture #1 \quad$\cdot$\quad Part V \quad$\cdot$\quad
     Extended notes}\\[3mm]
    {\Huge\bfseries\color{aimlprimary}\baselineskip=1.15\baselineskip #2\par}\vspace{4mm}
    {\color{aimlmuted}\large Dataset: #3}\\[3mm]
    {\color{aimlrule}\rule{\textwidth}{0.8pt}}
  \end{flushleft}
  \vspace{2mm}
  \begin{keybox}[title={Why these notes exist}]
    Lectures 19--22 sit outside G\'eron, the course's primary text. For those
    four lectures \emph{these notes are the primary source}, and they are
    examined on exactly the same terms as the textbook parts. Everything in
    the slide deck for this lecture is here, with the arguments written out at
    length rather than compressed onto a slide; the numbers are the ones the
    accompanying notebook prints, and every one of them is a measurement on
    the corpus named above rather than a figure quoted from a paper.
  \end{keybox}
  \vspace{1mm}
  {\small\tableofcontents}
  \vspace{2mm}
  \noindent{\color{aimlrule}\rule{\textwidth}{0.4pt}}
  \vspace{1mm}
}


\begin{document}
\notestitlefor{8}{I}{Dimensionality reduction and unsupervised learning}{Olivetti faces, 400 photographs}{Chapters 7--8}

\section{The first lecture with no target column}

Seven lectures, and every one was handed the answers: a price for every
district, a digit for every image, a survivor flag for every passenger, a cover
type for every plot. Today there is nothing. That changes what a metric can be,
what a baseline is, and how you find out you were wrong.

\begin{center}
\small
\begin{tabular}{p{0.42\textwidth}p{0.44\textwidth}}
\toprule
\textbf{Supervised} & \textbf{Unsupervised} \\
\midrule
You are given $(\mathbf{x}^{(i)}, y^{(i)})$ & you are given $\mathbf{x}^{(i)}$
                                              only \\
The metric compares prediction to truth & the metric compares the grouping to
                                          itself \\
The baseline is ``predict the majority'' & the baseline is ``group at
                                           random'' \\
You know when you are wrong & you find out when someone looks \\
\bottomrule
\end{tabular}
\end{center}

\begin{keybox}[title={The consequence for the whole hour}]
Every number computed today is a statement about \emph{geometry}, not about
people. Whether the geometry is the people is a separate question, and it is
answered by looking.
\end{keybox}

\subsection{The brief}

``We have a photographic archive. We do not know how many distinct people are
in it. Group the photographs by person, so that a human can name a group
instead of naming every photograph.'' Nobody remembers who is in it, an
annotator costs money per photograph, and we are funded for \emph{forty}
labels. That is the ordinary situation outside a benchmark: data is cheap,
labels are not.

400 photographs of 40 people, ten each, $64\times64$ pixels flattened to
$\mathbf{x}^{(i)} \in \mathbb{R}^{4096}$ with $x_j \in [0,1]$ --- 6.6 MB for
the whole corpus. Ten times more dimensions than points.

\begin{failbox}[title={The rule that keeps this honest}]
Olivetti ships with the identity of every photograph; the archive we are
pretending to hold does not. So $y$ is used in exactly two ways, both announced
where they happen: as the \emph{forty labels we paid for}, and as an
\emph{audit}, to find out what those forty could not have told us.

Any other use of $y$ is the same error as scoring on the test set, wearing
different clothes.
\end{failbox}

\subsection{Before anything is fitted: can this work at all?}

$k$-means groups points close in Euclidean distance, so one measurement decides
whether the hour is worth spending. It is an audit --- it uses the hidden
labels.

\begin{center}
\small
\begin{tabular}{lrrr}
\toprule
\textbf{Euclidean distance, raw pixels} & \textbf{Pairs} & \textbf{Mean} &
\textbf{Median} \\
\midrule
Same person      &  1{,}800 &  8.38 &  8.15 \\
Different people & 78{,}000 & 12.41 & 12.15 \\
\bottomrule
\end{tabular}
\end{center}

The gap is real and it is small: 6\% of same-person pairs sit further apart
than the \emph{median} different-person pair.

\begin{failbox}[title={Failure condition --- stated in advance}]
No method that groups by Euclidean distance in raw pixels will recover the
forty people cleanly. Not because the method is bad --- because the
representation puts a lamp and a nose on the same axis, with the same weight.

Moving the light source changes hundreds of pixel values by a large amount.
Changing the person changes a few dozen by a small one. Euclidean distance adds
them all up equally.

Measuring this costs one cell and tells you whether you are about to tune or
about to hope.
\end{failbox}

It is also not free to work there. One model-selection sweep --- $k$-means at
every $k$ from 2 to 60, ten restarts each, scored by silhouette --- costs 77
seconds for 590 runs of Lloyd's algorithm at $O(\text{iterations}\cdot mkn)$
and 2 more for 59 silhouette scores at $O(m^2 n)$. Both cost models are linear
in $n$. Two of the three factors are properties of the problem; the third is a
property of the \emph{representation}.

\section{Derivation: PCA via the SVD}

Two questions. Among all $d$-dimensional subspaces, which loses the least when
you project onto it, and how much does it lose? And if all you need is that
\emph{distances} survive, how small can $d$ be?

Write $\mathbf{X} \in \mathbb{R}^{m\times n}$ with one photograph per row,
$\tilde{\mathbf{X}} = \mathbf{X} - \mathbf{1}\bar{\mathbf{x}}^{\intercal}$ the
centred data, $\mathbf{W} \in \mathbb{R}^{n\times d}$ with
$\mathbf{W}^{\intercal}\mathbf{W} = \mathbf{I}_d$, and
$\mathbf{P} = \mathbf{W}\mathbf{W}^{\intercal}$ the orthogonal projector.

\begin{failbox}[title={Failure condition --- an SVD taken without centring}]
Without centring, the direction of largest $\sum_i\lVert\mathbf{w}^{\intercal}
\mathbf{x}^{(i)}\rVert^2$ is essentially $\bar{\mathbf{x}}$ itself, and the
first component describes nothing but ``this is a photograph of a face''.

Scikit-Learn's \code{PCA} centres for you. If you write the SVD by hand it does
not, and the result is a plausible-looking first component and no error
message. It is the commonest silent bug in a hand-rolled PCA.
\end{failbox}

\begin{derivbox}[title={Seven steps}]
\textbf{1. The objective.} Project each centred face and measure how far it had
to move:
\[ E(\mathbf{W}) = \sum_i \bigl\lVert\tilde{\mathbf{x}}^{(i)}
   - \mathbf{W}\mathbf{W}^{\intercal}\tilde{\mathbf{x}}^{(i)}\bigr\rVert_2^2
   = \bigl\lVert\tilde{\mathbf{X}}
   - \tilde{\mathbf{X}}\mathbf{W}\mathbf{W}^{\intercal}\bigr\rVert_F^2. \]

\textbf{2. Pythagoras.} $\mathbf{P}$ is an orthogonal projector, so kept and
lost are orthogonal, and since $\lVert\tilde{\mathbf{X}}\rVert_F^2$ does not
depend on $\mathbf{W}$,
\[ E(\mathbf{W}) = \lVert\tilde{\mathbf{X}}\rVert_F^2
                 - \lVert\tilde{\mathbf{X}}\mathbf{W}\rVert_F^2. \]

\textbf{3. As a trace.}
$\lVert\tilde{\mathbf{X}}\mathbf{W}\rVert_F^2
 = \operatorname{tr}(\mathbf{W}^{\intercal}\mathbf{C}\mathbf{W})$ with
$\mathbf{C} = \tilde{\mathbf{X}}^{\intercal}\tilde{\mathbf{X}}$.

\textbf{4. The SVD diagonalises it.} With
$\tilde{\mathbf{X}} = \mathbf{U}\boldsymbol\Sigma\mathbf{V}^{\intercal}$,
\[ \mathbf{C} = \mathbf{V}\boldsymbol\Sigma^2\mathbf{V}^{\intercal}, \]
so the columns of $\mathbf{V}$ are the eigenvectors of $\mathbf{C}$ with
eigenvalues $\sigma_j^2$. They are the principal components.

\textbf{5. Rotate the unknown.} Write $\mathbf{W} = \mathbf{V}\mathbf{A}$ with
$\mathbf{A}^{\intercal}\mathbf{A} = \mathbf{I}_d$; then
$\operatorname{tr}(\mathbf{W}^{\intercal}\mathbf{C}\mathbf{W}) =
 \sum_j \sigma_j^2 h_j$, where $h_j$ is the squared norm of row $j$ of
$\mathbf{A}$, constrained to $0 \le h_j \le 1$ with $\sum_j h_j = d$.

\textbf{6. A weighted sum with $d$ units to spend.} The $\sigma_j^2$ are
sorted, so put a full unit on each of the $d$ largest and nothing on the rest:
the maximum is $\sum_{j\le d}\sigma_j^2$, attained at $\mathbf{W} =
\mathbf{V}_d$.

\textbf{7. The error that is left.} Since
$\lVert\tilde{\mathbf{X}}\rVert_F^2 = \operatorname{tr}(\mathbf{C})
 = \sum_j\sigma_j^2$,
\[ \min_{\mathbf{W}} E(\mathbf{W}) = \sum_{j > d}\sigma_j^{2}. \]
\end{derivbox}

That is Eckart--Young--Mirsky in the Frobenius norm. ``PCA minimises
reconstruction error'' is a \emph{theorem}, not a slogan --- and it names the
price in advance: the tail of the spectrum, readable before you choose $d$.

Note also that step 2 says minimising error and maximising projected variance
are \emph{the same problem}, not two related ones. And $\mathbf{C}$ is the
matrix Lecture 2's normal equation had to invert, so step 4 says exactly when
it cannot be: when some $\sigma_j = 0$.

\begin{keybox}[title={A theorem you can check to machine precision}]
\begin{verbatim}
lhs = ((X_centred - Xd) ** 2).sum()   # the rank-d truncation's error
rhs = (S[d:] ** 2).sum()              # the tail of the spectrum
assert abs(lhs - rhs) / rhs < 1e-5
\end{verbatim}
Two sides computed by entirely different routes. When a theorem hands you an
identity, \emph{assert the identity}: it is the cheapest possible test that
your code is the theorem.
\end{keybox}

\subsection{Choosing $d$ without guessing}

By step 7 the fraction of squared error removed by keeping $d$ components is
exactly $\sum_{j\le d} r_j$, where $r_j = \sigma_j^2/\sum_i\sigma_i^2$. The
explained variance ratio is not a heuristic; it is the theorem, normalised.

The first component alone holds 23.8\%, ten hold 65.6\%, fifty hold 87.4\%.
Cumulating to a target: 66 components reach 90\%, \textbf{123 reach 95\%}, 260
reach 99\%. The rest of this lecture runs at 123 --- a factor of 33 fewer
numbers per photograph.

A principal component is a vector in $\mathbb{R}^{4096}$, so reshaping it to
$64\times64$ makes it an image; on faces these are called \term{eigenfaces}.
That is not a property of faces but of PCA: the components are directions in
the input space, so whatever an input looks like, a component looks like too.
And looking at the first few, they are \emph{lighting} --- exactly what the
distance histogram predicted would dominate.

\subsection{The second question: what if you only need distances?}

PCA looks at the data and spends an SVD finding the best subspace. The
alternative is to not look at the data at all: draw
$\mathbf{R} \in \mathbb{R}^{n\times d}$ with independent Gaussian entries and
set $\mathbf{Z} = \frac{1}{\sqrt d}\mathbf{X}\mathbf{R}$.

\begin{derivbox}[title={The Johnson--Lindenstrauss lemma}]
For any $0 < \varepsilon < 1$ and any set of $m$ points in $\mathbb{R}^n$,
there is a linear map $f$ into $\mathbb{R}^d$ with
\[ d \;\ge\; \frac{4\log m}{\varepsilon^{2}/2 - \varepsilon^{3}/3} \]
such that for every pair $\mathbf{u},\mathbf{v}$,
\[ (1-\varepsilon)\lVert\mathbf{u}-\mathbf{v}\rVert^2
   \le \lVert f(\mathbf{u})-f(\mathbf{v})\rVert^2
   \le (1+\varepsilon)\lVert\mathbf{u}-\mathbf{v}\rVert^2. \]
The tolerance is on \emph{squared} distances. Quoting $\varepsilon$ against
unsquared ones understates the distortion by about a factor of two.
\end{derivbox}

\begin{keybox}[title={Read the formula again. What is missing?}]
In it: $m$, the number of points, logarithmically; and $\varepsilon$, the
tolerance. \emph{Not} in it: $n$, the dimension you are starting from.

400 points need the same target dimension whether they live in 4{,}096
dimensions or in four million.
\end{keybox}

\begin{center}
\small
\begin{tabular}{lrr}
\toprule
$\varepsilon$ & \textbf{400 points} & \textbf{One million points} \\
\midrule
0.1 & 5{,}135 & 11{,}841 \\
0.2 & 1{,}382 &  3{,}188 \\
0.3 &   665 &  1{,}535 \\
0.5 &   287 &    663 \\
\bottomrule
\end{tabular}
\end{center}

Two and a half thousand times as many points costs about twice the dimension.
That is the $\log m$, and it is why the theorem is interesting at all.

But be honest about the top-left cell: 5{,}135 dimensions to hold every
pairwise distance among 400 faces to within 10\% is \emph{more than the 4{,}096
we started with}. The bound is a worst case over all point sets --- it has to
work for the nastiest configuration anyone can construct, and a photographic
archive is not that. So it is sufficient, never necessary, and the reason to
teach it is not the constant. Measured on our own distances at the bound's own
$d = 1{,}382$, the worst of the 79{,}800 pairs is distorted by 0.17: inside the
0.2 promised, and not by much.

\section{Reducing the dimension}

\begin{center}
\small
\begin{tabular}{llr r}
\toprule
\textbf{Method} & \textbf{What it does} & \textbf{Fit time, s} &
\textbf{Held-out error} \\
\midrule
PCA, full SVD    & the complete SVD, then truncate & 0.074 & 0.00240 \\
PCA, randomised  & approximates the top $d$        & 0.081 & 0.00241 \\
Incremental PCA  & one batch at a time             & 0.142 & 0.00240 \\
Random projection& never looks at the data         & 0.006 & 0.32093 \\
\bottomrule
\end{tabular}
\end{center}

Three PCA variants land on the same error at different prices. Random
projection is essentially free and pays with a reconstruction error 133 times
larger --- a random 123-dimensional subspace of $\mathbb{R}^{4096}$ keeps only
about $123/4096 = 3.0\%$ of a generic vector's squared norm.

\begin{keybox}[title={That is not a bug in random projection}]
PCA minimises exactly this column by construction. Johnson--Lindenstrauss
promises something else entirely --- that \emph{relative} distances survive.
Neither theorem is about the other quantity.

So the question is never which is better. It is whether you need the
reconstruction or only the distances.
\end{keybox}

PCA is the wrong tool under three conditions. \emph{Variance is not signal}: if
the largest variance is a nuisance --- a lamp --- PCA keeps the nuisance first.
\emph{The structure is not linear}: a subspace is flat. \emph{The features are
not commensurable}: PCA is not scale-free, so standardise first unless, as
here, every feature is already the same kind of number. The first condition is
the one on this data, and no amount of compression repairs it.

\subsection{One procedural condition, because it is invisible}

A reducer is an estimator: it is \emph{fitted}. Fit it before the split and the
held-out faces have helped choose the subspace they are then judged in. The 280
training faces span at most a 280-dimensional subspace of $\mathbb{R}^{4096}$,
so adding 120 more points genuinely changes which directions survive.

\begin{center}
\small
\begin{tabular}{lrrl}
\toprule
\textbf{$d = 123$, 20 stratified splits} & \textbf{Fitted on train} &
\textbf{Fitted on everything} & \textbf{Leaky wins} \\
\midrule
Reconstruction error, held out & 0.00243 & 0.00096 & 20 of 20 \\
Accuracy, held out             & 97.4\%  & 97.4\%  & 3 of 20 \\
\bottomrule
\end{tabular}
\end{center}

\begin{failbox}[title={Failure condition --- the decision rule, and why it is still procedural}]
The leak makes the reconstruction error look 61\% better and moves the accuracy
by nothing measurable.

Fitting an unsupervised step on everything costs you exactly when \emph{that
step's own objective is the number you report}. Reconstruction error is that
objective; downstream accuracy is not --- and you cannot tell which row you are
in without doing the split. So put it in a \code{Pipeline}, which refits inside
every fold and makes the leak structurally impossible rather than merely
avoided.
\end{failbox}

\section{Clustering, and how to choose $k$}

$k$-means makes four assumptions whether you asked or not: the partition is a
Voronoi diagram, so every cluster is \emph{convex}; the only quantity is
distance to a centre, so clusters are \emph{spherical and equal-sized}; every
point belongs to exactly one cluster, so there is no ``none of these''; and
\emph{$k$ is given}. Three of the four are false here.

\begin{center}
\small
\begin{tabular}{llc}
\toprule
\textbf{Metric} & \textbf{Needs} & \textbf{Available here?} \\
\midrule
RMSE, accuracy, $F_1$, ROC AUC & a truth per instance & no \\
Inertia   & nothing & yes \\
Silhouette & nothing & yes \\
Adjusted Rand index & labels --- a subset is enough & on the forty \\
\bottomrule
\end{tabular}
\end{center}

Two of these can choose a model. One can be shown to a stakeholder. They are
not the same one, and keeping those jobs apart is this course's standing rule
applied where there is no test set at all.

\subsection{Why inertia cannot choose $k$}

\[ J(k) = \sum_i \bigl\lVert\mathbf{x}^{(i)}
          - \boldsymbol\mu_{c(i)}\bigr\rVert_2^2 \]
is what $k$-means minimises --- and it is non-increasing in $k$, with
$J(m) = 0$ exactly. Refining a partition cannot increase the sum of squared
distances to the assigned centre. \emph{There is no interior minimum to find,
at any $k$, on any dataset}, so ``minimise inertia'' selects one cluster per
photograph. The elbow rule exists to work around that, and on our curve the
kneedle rule picks $k = 19$ where the archive has 40 --- with 46\% of the
inertia at $k=2$ still present at $k=40$.

\begin{keybox}[title={The elbow is not a criterion}]
It is a picture of a criterion that does not exist, and no better heuristic
repairs it, because the information is not in $J$.
\end{keybox}

The silhouette does have an interior maximum. For a point $i$, with $a$ its
mean distance to its own cluster and $b$ to the nearest other,
\[ s_i = \frac{b-a}{\max(a,b)} \in [-1,1], \]
penalised both for splitting a group and for merging two. It costs $O(m^2n)$
--- every pair, over every feature, and both factors matter.

\subsection{Before reading any silhouette: what does nothing score?}

\begin{center}
\small
\begin{tabular}{lrr}
\toprule
\textbf{Random assignment} & \textbf{Silhouette} & \textbf{ARI on the forty} \\
\midrule
$k = 10$ & $-0.0422 \pm 0.0045$ & $-0.0040 \pm 0.0205$ \\
$k = 40$ & $-0.1288 \pm 0.0107$ & $-0.0089 \pm 0.0244$ \\
$k = 60$ & $-0.1805 \pm 0.0124$ & $-0.0075 \pm 0.0261$ \\
\bottomrule
\end{tabular}
\end{center}

Both anchors sit at zero, and now you know the width of the noise around zero.
A silhouette of 0.02 is not a discovery.

The \term{adjusted Rand index} compares two groupings pair by pair and
subtracts what chance would give, so random scores 0 in expectation at any $k$
and perfect scores 1. We may compute it on the forty labelled photographs
only --- 780 pairs, of which just 22 are the same person --- so treat it as a
sanity check with a wide interval, never as a ranking.

\begin{center}
\small
\begin{tabular}{lrrr}
\toprule
& $k$ & \textbf{Silhouette} & \textbf{ARI on the forty} \\
\midrule
Random assignment      & 40 & $-0.1288$ & $-0.0089$ \\
The silhouette's choice & 59 & 0.1707 & 0.3954 \\
The true $k$ --- an audit & 40 & 0.1465 & 0.2978 \\
\bottomrule
\end{tabular}
\end{center}

Well above the anchor, so the clustering is doing something. It chose 59 groups
for 40 people, and the silhouette at the truth is \emph{lower} than at its own
choice.

\subsection{The mean hides the shape}

A silhouette diagram draws one knife per cluster, width for size and length for
quality. A knife entirely left of the mean marks a cluster worse than average;
knives of very unequal width mean one group has swallowed several; a knife
extending into negative $s_i$ marks points closer to another cluster than their
own; and \emph{every knife short} means the clusters overlap and the geometry
has no clean gaps. Ours is the last, at every $k$ we tried --- which is the
diagnosis the mean of 0.1707 does not contain.

\begin{center}
\small
\begin{tabular}{lr}
\toprule
\textbf{At $k = 40$, against the hidden identities} & \textbf{Value} \\
\midrule
Cleanest cluster: 10 photographs, purity & 1.00 \\
Worst cluster: 24 photographs of 10 people, purity & 0.21 \\
Mean purity across the forty clusters & 0.74 \\
Clusters holding exactly one person & 16 \\
Clusters with exactly ten members & 7 \\
Cluster sizes, smallest to largest & 4 to 25 \\
\bottomrule
\end{tabular}
\end{center}

Every true group has exactly ten members. The failure condition stated at the
start of the lecture is now a measurement --- and the worst cluster is not a
person at all, but a lighting condition: several different people lit and posed
the same way.

\subsection{The score that chose the model also reports it}

\begin{failbox}[title={Failure condition --- a bias that needs no labels to exist}]
Sweep $k$, keep the best silhouette, print it, and the number that chose the
model is the number that reports it. Measured: choose $k$ on a random half of
the corpus and score the chosen model on the other half, twenty times.

\begin{center}
\small
\begin{tabular}{lrr}
\toprule
\textbf{Silhouette, 20 random halves} & \textbf{Mean} & \textbf{sd} \\
\midrule
Reported --- chose and scored on the same half & 0.1777 & 0.0123 \\
Held out --- had no vote & 0.0878 & 0.0154 \\
Optimism & 0.0899 & 0.0229 \\
\bottomrule
\end{tabular}
\end{center}
Half the reported score, and worse in 20 splits out of 20.
\end{failbox}

\section{Putting the two halves together}

The clustering code is unchanged --- same grid, same seed, same
\code{n\_init}. Only what $k$-means is looking at changed.

\begin{center}
\small
\begin{tabular}{lrrrr}
\toprule
\textbf{Dimensions} & \textbf{Sweep, s} & \textbf{Best $k$} &
\textbf{Silhouette} & \textbf{ARI, all 400} \\
\midrule
4{,}096 --- raw pixels & 78 & 59 & 0.171 & 0.537 \\
20  & 1.9 & 58 & 0.286 & 0.540 \\
50  & 1.7 & 59 & 0.256 & 0.493 \\
100 & 1.9 & 60 & 0.218 & 0.547 \\
123 --- 95\% of the variance & 2.1 & 60 & 0.205 & 0.524 \\
399 & 5.8 & 59 & 0.171 & 0.537 \\
\bottomrule
\end{tabular}
\end{center}

The sweep is 37 times faster at 123 dimensions and the agreement with the true
identities does not fall. The ARI column is there because a speed-up that
quietly costs quality is not a speed-up. The seconds are a property of one
machine at two threads; the \emph{shape} --- an order-of-magnitude saving at no
cost in ARI --- is a property of the method.

\begin{failbox}[title={Failure condition --- careful: the silhouette went up}]
That is not evidence the clustering improved. The silhouette is computed in
whatever space you hand it, and compressing the data changes every distance, so
the criterion is \emph{not comparable across the rows of that table}.

Which is exactly why the ARI column is there: it is measured against the
identities and does not depend on the representation. Never compare two scores
measured on different evaluation sets --- and a different representation is a
different evaluation set.
\end{failbox}

\section{Beyond $k$-means}

Compression fixed the clock. It did not fix the geometry.

\begin{center}
\small
\begin{tabular}{lll}
\toprule
& \textbf{DBSCAN} & \textbf{Gaussian mixture} \\
\midrule
Cluster shape & any connected dense region & an ellipsoid \\
Number of clusters & found, not given & given \\
Every point assigned? & no --- there is a noise label & softly, with a
                                                       probability \\
Main parameter & \code{eps}, a distance & \code{n\_components} \\
\bottomrule
\end{tabular}
\end{center}

\begin{keybox}[title={Reporting a method that did not work}]
DBSCAN needs one density scale that separates clusters from gaps. Best over the
whole grid: ARI 0.133 at \code{eps=6.0}, with 41 clusters and 97 of the 400
faces called noise; it never finds more than 50 clusters and never 40.

Forty groups of ten faces, all lit by the same rig, have no such scale --- the
gap between two people is about the gap between two photographs of one person.
This is a real result and it belongs in the report. A method that fails for a
stateable reason is more informative than a method that was never tried.
\end{keybox}

A Gaussian mixture models the data as $k$ Gaussians fitted by
expectation--maximisation, so clusters may be elongated and of different sizes,
assignment is soft, and --- being a density model --- it also scores how likely
a new point is.

\begin{failbox}[title={Why this needed the first half of the lecture}]
A full covariance in $n$ dimensions has $n(n+1)/2$ free parameters.

\begin{center}
\small
\begin{tabular}{lr}
\toprule
\textbf{Setting} & \textbf{Free parameters} \\
\midrule
One full covariance in 4{,}096 dimensions & 8{,}390{,}656 \\
Forty of them & 335{,}626{,}240 \\
Forty diagonal components in 123 dimensions --- the whole model & 9{,}880 \\
\bottomrule
\end{tabular}
\end{center}

Estimated from 400 photographs. The first is not slow --- it is
\emph{undefined}: every covariance estimate is singular. Arithmetic settles
that in one line, where an afternoon of waiting does not.
\end{failbox}

BIC picks 5 components while the ARI peaks at 55. Both penalise complexity, but
they answer a different question: BIC asks which model best explains this data
\emph{as a density}, and we are asking which grouping matches the people. Those
coincide only when the people really are Gaussian blobs, and here they are not.

\subsection{Anomaly detection, two ways}

Plant twelve corrupted images --- four rotated, four dimmed, four
double-exposed --- and count how many each detector puts in its top twelve.
Reconstruction error caught \textbf{12 of 12}; the mixture density caught 6.

A rotated or double-exposed face is outside the subspace the eigenfaces span,
so its reconstruction error is large: median 0.0017 against 0.0002 for a
genuine face. The mixture density is estimated \emph{inside} that subspace, so
a corruption projecting onto an ordinary-looking point is invisible to it. Use
reconstruction error when the anomaly is a different \emph{kind} of object; use
density when it is an ordinary object in an unusual \emph{place}.

\section{Spending the forty labels}

Where they are spent turns out to matter more than what is done with them. Same
classifier, same held-out faces, twenty splits; the only thing that varies is
which forty were labelled.

\begin{center}
\small
\begin{tabular}{lrr}
\toprule
\textbf{Labels used} & \textbf{Held-out accuracy} & \textbf{sd over 20 splits} \\
\midrule
40 at random                        & 48.1\% & 3.8\% \\
40, one per cluster                 & 60.5\% & 4.9\% \\
Propagated to the whole cluster     & 55.7\% & 4.7\% \\
Propagated to the closest 75\%      & 55.1\% & 4.8\% \\
All 280 true labels --- the ceiling & 97.4\% & 1.6\% \\
\bottomrule
\end{tabular}
\end{center}

Twenty splits, because the differences between the middle three rows are of the
same order as their spread. \emph{Propagation buys 280 labels from 40 and does
worse than the 40 alone}, because a wrong propagated label is worse than no
label. Report that; do not bury it.

\subsection{The diagnosis was right, and the obvious repair still fails}

\begin{center}
\small
\begin{tabular}{lrrr}
\toprule
\textbf{Representation} & \textbf{Best $k$} & \textbf{Silhouette} &
\textbf{ARI, all 400} \\
\midrule
PCA to 123 dimensions & 60 & 0.205 & 0.524 \\
PCA 123, first 3 components discarded & 60 & 0.184 & 0.508 \\
Unit-norm each image, then PCA 123 & 58 & 0.189 & 0.464 \\
\bottomrule
\end{tabular}
\end{center}

Neither helps. The leading components carry pose and face shape as well as
illumination, so removing them throws away signal with the nuisance. The remedy
that works is a representation learned \emph{for the task} rather than for
variance --- which is what Part VI builds, and why a pretrained encoder beats a
hand-chosen projection.

Note also what chose $k$ there: sweeping $k$ and keeping the best ARI would
make both repairs look like gains, because ARI is computed from the labels this
whole lecture is pretending not to have.

\section{The notebook}

\begin{trybox}[title={What to change}]
Change \code{n\_components=0.95} to \code{0.99}, which takes $d$ from 123 to
260. The sweep slows by roughly the ratio of the two and the ARI barely moves
--- the whole trade-off in one edit.
\end{trybox}

\section{Exercises}

\begin{enumerate}[leftmargin=1.6em,itemsep=4pt]
  \item PCA is derived by maximising $\lVert XW\rVert_F^2$ subject to
        $W^{\mathsf T}W = I$. State why the data must be centred first, and
        what the first component becomes if it is not. \hfill \scope{[5]}
  \item The Johnson--Lindenstrauss bound, evaluated for 400 images at
        $\varepsilon = 0.1$, asks for more dimensions than the data has. What
        do you conclude about the bound, and about the method?
        \hfill \scope{[5]}
  \item A pipeline fits PCA on the whole dataset and then splits. Name what
        leaked, and say whether the leak grows or shrinks as the corpus grows.
        \hfill \scope{[4]}
  \item Inertia always falls as $k$ rises, so it cannot choose $k$. State what
        the silhouette adds, and one situation where it also fails.
        \hfill \scope{[5]}
  \item You have 4{,}096-dimensional face vectors and are told two faces are
        ``close''. Explain why that alone does not mean they are the same
        person. \hfill \scope{[4]}
\end{enumerate}

\section*{Summary}
\addcontentsline{toc}{section}{Summary}

PCA is the SVD of the centred data: the principal subspace minimises squared
reconstruction error exactly, and the error it leaves is
$\sum_{j>d}\sigma_j^2$, readable before you choose $d$. Minimising that error
and maximising the projected variance are the same problem by Pythagoras, and
the explained variance ratio is that theorem normalised --- 123 components hold
95\% of 4{,}096 pixels. Johnson--Lindenstrauss bounds the target dimension by
the number of points and the tolerance, never by the dimension you started in,
and it is a loose worst case, so measure the distortion you actually get.

Inertia cannot choose $k$: it is monotone in $k$ and zero at $k = m$. The
silhouette can, at $O(m^2n)$, must be read against a random anchor, and is
optimistic by about half when it both chooses and reports. Compressing to 123
dimensions made the sweep 37 times faster at no cost in ARI --- but the
silhouette rose, and that number is not comparable across representations.

The lecture opened by measuring that pixel distance is partly identity and
partly lighting, and closed by confirming it: mean cluster purity 0.74, the
worst cluster a lighting condition rather than a person, DBSCAN failing for a
stateable reason, and two one-line repairs that both made things worse. And
forty labels spent on cluster representatives were worth a quarter more
accuracy than forty spent at random --- while propagating them to 280 was worse
than not propagating at all.

\end{document}
